\begin{recipe}{Döner kebab uit de oven}
\kicker{Hoofdgerecht · vlees · Turks}
\index[register]{Döner kebab}
\index[register]{Gehakt}
\index[register]{Lamsgehakt}

\meta{50 minuten \dvd Voor 5 personen}

\lettrine{D}{öner kebab is een Turkse klassieker, maar thuis zonder draaiende spit even lekker te maken. Dit recept rolt het gehakt flinterdun uit tussen bakpapier en bakt het in de oven tot een krokante buitenkant. Het resultaat is vochtiger en smaakvoller dan je verwacht.}

\blockrule{Ingrediënten}
\begin{ingredients}
  \ing{500~g}{lamsgehakt}
  \ing{500~g}{rundergehakt}
  \ing{2}{uien}
  \ing{4 tenen}{knoflook}
  \ing{4 el}{Griekse yoghurt}
  \ing{4 el}{olijfolie}
  \ing{2 tl}{zout}
  \ing{3 tl}{paprikapoeder}
  \ing{2 tl}{gedroogde oregano}
  \ing{2 tl}{komijnpoeder}
  \ing{3 tl}{Pul Biber}
  \ing{1 tl}{zwarte peper}
\end{ingredients}

\blockrule{Bereiding}
\tip{Het kneden van het gehakt minimaal 3 tot 5 minuten is essentieel: het vlees moet glad en kleverig worden, bijna als paté, zodat het straks goed aan elkaar blijft zitten.}
\begin{steps}
  \item Blend de uien tot een gladde puree. Rasp de knoflook fijn tot pasta.
  \item Doe het gehakt, yoghurt, olijfolie, uienpuree, knoflookpasta en kruiden in een grote kom. Kneed dit met je handen intensief 3 tot 5 minuten, of gebruik een keukenmachine met het mes-opzetstuk totdat het vlees glad en kleverig is.
  \item Verdeel de vleesmassa in 4 gelijke porties. Leg een groot vel bakpapier op het aanrecht, plaats één portie vlees in het midden en leg er een tweede vel bakpapier bovenop.
  \item Rol het vlees met een deegroller flinterdun uit tot 2 tot 3 millimeter dik.
  \item Haal het bovenste vel papier eraf. Pak het onderste vel aan de lange zijde vast en rol het vleesblad strak op tot een lange, compacte cilinder, zonder het papier erin mee te rollen. Wikkel de rol in het papier en draai de zijkanten dicht als een snoepje. Herhaal voor alle 4 porties.
  \item Verwarm de oven voor op 200°C. Leg de vier ingepakte rollen op een bakplaat en bak ze 20 minuten.
  \item Haal de bakplaat uit de oven. Giet het vet en vocht voorzichtig weg.
  \item Haal de rollen uit het bakpapier en leg de kale rollen terug op de droge bakplaat. Bak nog 15 minuten totdat de buitenkant goudbruin is.
  \item Laat het vlees 5 minuten rusten. Snijd het in flinterdunne plakjes met je scherpste mes.
  \item Voor extra krokante randjes: verwarm de ovengrill of een droge koekenpan op hoog vuur en bak de gesneden plakjes 3 tot 4 minuten tot de randjes donkerbruin worden.
  \item Serveer in getoast pita of platbrood met knoflooksaus en een frisse uien-sumaksalade.
\end{steps}
\end{recipe}
